\documentclass[a4paper,landscape,headrule,footrule]{foils}

%%
%%% macros for Theories of Grammar
%%%
\usepackage{polyglossia}
\setdefaultlanguage{english}
%\setmainfont{TeX Gyre Pagella}


\newcommand{\logo}{~}
\newcommand{\header}[3]{%
\title{\vspace*{-2ex} \large JMORF/MORS --- Morpho-Syntax
\\[2ex] \Large  \emp{#2} \\ \emp{#3}}
\author{\blu{Francis Bond}   \\ 
\normalsize  \textbf{Palacký University}\\
\normalsize  \url{https://fcbond.github.io/}\\
\normalsize  \texttt{bond@ieee.org}}
\MyLogo{JMORF (2026)}
\renewcommand{\logo}{#2}
\hypersetup{
   pdfinfo={
     Author={Francis Bond},
     Title={#1: #2},
     Subject={JMORF/MORS: Morpho-Syntax},
     Keywords={Syntax, Morphology, HPSG, Unification, Constructions},
     License={CC BY 4.0}
   }
 }

\date{#1
  \\ Location: KB-1.11}
}
\usepackage{hyperref}
\hypersetup{
     colorlinks,
     linkcolor={blue!50!black},
     citecolor={red!50!black},
     urlcolor={blue!80!black}
   }



\usepackage{xcolor}
\usepackage{graphicx}
\newcommand{\blu}[1]{\textcolor{blue}{#1}}
\newcommand{\grn}[1]{\textcolor{green}{#1}}
\newcommand{\hide}[1]{\textcolor{white}{#1}}
\newcommand{\emp}[1]{\textcolor{red}{#1}}
\newcommand{\txx}[1]{\textbf{\textcolor{blue}{#1}}}
\newcommand{\lex}[1]{\textbf{\mtcitestyle{#1}}}


% \renewcommand{\labelitemi}{\textcolor{violet}{‣}}
% \renewcommand{\labelitemii}{\textcolor{purple}{‣}}

\newcommand{\subhead}[1]{\noindent\textbf{#1}\\[5mm]}

\newcommand{\Bad}{\emp{\raisebox{0.15ex}{\ensuremath{\mathbf{\otimes}}}}}
\newcommand{\bad}[1]{*\eng{#1}}

\newcommand{\com}[1]{\hfill (\emp{#1})}%

\usepackage{relsize,xspace}
\newcommand{\into}{\ensuremath{\rightarrow}\xspace}
\newcommand{\tot}{\ensuremath{\leftrightarrow}\xspace}
\usepackage{url}
\newcommand{\lurl}[1]{\MyLogo{\url{#1}}}

\usepackage{langsci-gb4e}
\let\eachwordone=\slshape
%%% langsci-gb4e defines \thexnumi as \@roman\@xsi{xnumi}, which only works
%%% inside an exe environment; some decks use xlisti on its own.  Restore
%%% the plain gb4e definition so that keeps working.
\makeatletter
\def\thexnumi{\@xsi{xnumi}}
\makeatother
\newcommand{\lx}[1]{\textbf{\mtciteform{#1}}}
\newcommand{\ix}{\ex\slshape}



\newcommand{\ent}{\ensuremath{\Rightarrow}\xspace}
\newcommand{\ngv}{\ensuremath{\not\Rightarrow}\xspace}
%\usepackage{times}
%\usepackage{nttfoilhead}
\newcommand{\myslide}[1]{\foilhead[-25mm]{\raisebox{12mm}[0mm]{\emp{#1}}}\MyLogo{\logo}}
\newcommand{\myslider}[1]{\rotatefoilhead[-25mm]{\raisebox{12mm}[0mm]{\emp{#1}}}}
%\newcommand{\myslider}[1]{\rotatefoilhead{\raisebox{-8mm}{\emp{#1}}}}

\newcommand{\section}[1]{\myslide{}{\begin{center}\Huge \emp{#1}\end{center}}}



\usepackage[lyons,j,e,k]{mtg2e}
%\renewcommand{\mtcitestyle}[1]{\textcolor{-red!75!green!50}{\textsl{#1}}}
\renewcommand{\mtcitestyle}[1]{\textcolor{teal}{\textsl{#1}}}
\newcommand{\cs}{\mtciteform}
\newcommand{\iz}[1]{\texttt{\textup{#1}}}
\newcommand{\gm}{\textsc}
\usepackage[normalem]{ulem}
\newcommand{\ul}{\uline}
\newcommand{\ull}{\uuline}
\newcommand{\wl}{\uwave}
\newcommand{\vs}{\ensuremath{\Leftrightarrow}~}
%%%
%%% Bibliography
%%%
\usepackage{natbib}
%\usepackage{url}
\usepackage{bibentry}



\usepackage[utf8]{inputenc}
%%% trees
\usepackage{multicol}
%%% pst-node: one example in lec-12-control draws a connector between
%%% two points with \rnode/\ncbar.  This used to come in with rtrees.
\usepackage{pst-node}
\usepackage{forest}
\useforestlibrary{linguistics}
\forestapplylibrarydefaults{linguistics}

%%% \makebox's optional arguments are read as AVM delimiters by
%%% langsci-avm, so hide them behind a macro with no visible brackets.
\newcommand{\lbox}[1]{\makebox[0mm][l]{#1}}

%%% Tree styles.
%%%   words  -- trees whose leaves are actual words (italicised, aligned)
%%%   cat    -- trees whose leaves are categories or feature structures
\forestset{
  words/.style={for tree={s sep=7mm, l sep=5mm,
                          if n children=0{font=\itshape, tier=terminal}{}}},
  cat/.style={for tree={s sep=7mm, l sep=5mm}},
}
\newcommand{\wtree}[1]{\begin{forest} words #1 \end{forest}}
\newcommand{\ctree}[1]{\begin{forest} cat #1 \end{forest}}
%%% the rule being applied at this step of a derivation
\newcommand{\step}[1]{\begin{center}\emp{#1}\end{center}}

%%% AVMs (Language Science Press)
\usepackage{langsci-avm}

%%% Macros the decks use, carried over from headx.tex
\newcommand{\blank}{\rule{3em}{1pt}\xspace}
\newcommand{\ft}[1]{\textsc{#1}}
\newcommand{\idn}{\mtciteform}    %%% Indonesian
\newcommand{\prd}[1]{\textbf{#1}}
\newcommand{\typ}[1]{\textit{#1}}
%%% \val: provide before renewing, in case a package defines it.
\providecommand{\val}[1]{}
\renewcommand{\val}[1]{\textit{#1}}
%%% Manning's \avmbox, still used in tree labels and in maths where a
%%% bare langsci \tag would not work.  \avm{\tag{...}} gives the same
%%% boxed tag as the ones inside feature structures.
\newcommand{\avmbox}[1]{\mbox{\avm{\tag{#1}}}}

%%% Ring a value in a dotted oval.  These used pstricks \ovalnode,
%%% which came in with rtrees; tikz (via forest) does the same job.
\newcommand{\OV}[1]{\tikz[baseline]{\node[draw,dotted,red,ellipse,
          inner sep=1pt,outer sep=0pt]{#1};}}
\newcommand{\OVB}[1]{\tikz[baseline]{\node[draw,dotted,blue,ellipse,
          inner sep=1pt,outer sep=0pt]{#1};}}

%%% From CSLI book
\newcommand{\mc}{\multicolumn}
\newcommand{\HD}{\textbf{H}\xspace}
\newcommand{\el}{< >}       %%% empty list; langsci-avm needs the space

\newcommand{\A}{\noindent\textbf{A}: }
\newcommand{\Q}{\noindent\textbf{Q}: }

%%% Question and answer on two pages, written once:
%%%   \qaslide{Title}{... question ... \answer{... answer ...}}
%%% prints the slide twice, first without the \answer parts and then
%%% with them, so the answer appears "on click".  The example counter
%%% is reset so both pages show the same example numbers.
\newif\ifshowanswer
\newcommand{\answer}[1]{\ifshowanswer#1\fi}
\newcounter{qaxnumi}
\newcommand{\qaslide}[2]{%
  \setcounter{qaxnumi}{\value{xnumi}}%
  \showanswerfalse\myslide{#1}#2%
  \setcounter{xnumi}{\value{qaxnumi}}%
  \showanswertrue\myslide{#1}#2%
  \showanswerfalse}
%\newcommand{\C}{\noindent\textbf{C}: }



\begin{document}

\header{Lecture 6}{Binding Theory}{}
\maketitle


\myslide{Overview}
\MyLogo{\citet{sag:wasow:bender:2003} --- Chapter 7}
\begin{itemize}
\item Review: what we are trying to do; last week's semantics
\item The informal binding theory from Chapter 1
\item Formalizing binding theory: \ft{arg-st}, the ARP and the Binding Principles
\item Binding and PPs
\item Imperatives
\item Review exercise: semantic identity tables
\item Problem Set 6: possessives, imperative `subjects' (SWB 6.3, 6.4, 7.2)
\end{itemize}



\section{Review}

\myslide{What We’re Trying To Do}
\begin{itemize}
\item Objectives
\begin{itemize}
\item Develop a theory of knowledge of language
\item Represent linguistic information explicitly enough to
distinguish well-formed from ill-formed expressions
\item Assign plausible semantic representations
\item Be parsimonious, capturing linguistically significant
generalizations
\end{itemize}
\item Why formalize?
\begin{itemize}
\item To formulate testable predictions
\item To check for consistency
\item To make it possible to get a computer to do it for us
\end{itemize}
\end{itemize}

\myslide{How We Construct Sentences}
\begin{itemize}
\item The components of our grammar
\begin{itemize}
\item Grammar rules
\item Lexical entries
\item Principles
\item Type hierarchy (very preliminary, so far)
\item Initial symbol (S, for now) \\[2ex]
  \begin{small}
   \avm{S =  [ head & \textit{verb}\\
    val & [ spr & <  > \\
    comps & < > ] ]
}
\end{small}   

\end{itemize}
\item We combine constraints from these components.
%\item Q: What says we have to combine them?
% \item More in \url{hpsg.stanford.edu/book/slides/Ch6a.pdf}, 
% \\ \url{hpsg.stanford.edu/book/slides/Ch6b.pdf}
\end{itemize}


\myslide{Review of Semantics}
\begin{itemize}
\item Which aspects of semantics we’ll tackle
\item Semantics Principles
\item Building semantics of phrases
\item Modification, coordination
\item Structural ambiguity
\end{itemize}


\myslide{Our Slice of a World of Meanings}

Aspects of meaning we \emp{won’t} account for (in this course)
\begin{itemize}
\item Pragmatics
\item Fine-grained lexical semantics
  \\ The meaning of \eng{life} is
  \begin{itemize}\addtolength{\itemsep}{2ex}
  \item \iz{life}  
or 
\textsf{life}$'$
or
  \avm{%\avmvalfont{\it}
    [ reln & life \\ inst & i ]
}
\item Not like WordNet: 
\iz{life$_1$}  $\subset$ \iz{being$_1$} $\subset$ \iz{state$_1$} \ldots

 % -- (a characteristic state or mode of living; "social life"; "city life"; "real life")
 %       => being, beingness, existence -- (the state or fact of existing; "a point of view gradually coming into being"; "laws in existence for centuries")
 %           => state -- (the way something is with respect to its main attributes; "the current state of knowledge"; "his state of health"; "in a weak financial state")
 %               => attribute -- (an abstraction belonging to or characteristic of an entity)
 %                   => abstraction, abstract entity -- (a general concept formed by extracting common features from specific examples)
 %                       => entity -- (that which is perceived or known or inferred to have its own distinct existence (living or nonliving))
 \end{itemize}
\item Quantification \hfill (covered lightly in the book)
\item Tense, mood, aspect \hfill (covered in the book)
 \end{itemize}


% RELN life
% INST i

\myslide{Our Slice of a World of Meanings}
\noindent\avm{[mode &  prop \\
 index &  $s$\\
 restr & < [ RELN & {\bf save}\\
               SIT & $s$\\
              SAVER &  $i$\\
              SAVED &  $j$ ] \ , 
              [RELN & {\bf name}\\ 
              NAME &  Chris\\
              NAMED & \ $i$]\ , 
              [RELN & {\bf name}\\
              NAME &  Pat\\
              NAMED & \ $j$] > ]
} 
\begin{quote}
  “\ldots the linguistic meaning of \eng{Chris saved Pat} is a proposition
  that will be true just in case there is an actual situation that
  involves the saving of someone named Pat by someone named Chris.”
  \begin{flushright}
    \citep[p.~140]{sag:wasow:bender:2003}
  \end{flushright}
\end{quote}

\myslide{Semantics in Constraint-Based Grammar}
\begin{itemize}
\item Constraints as (generalized) truth conditions
  \begin{itemize}
\item \txx{proposition}: what must be the case for a proposition to be true
\item \txx{directive}: what must happen for a directive to be fulfilled
\item \txx{question}: the kind of situation the asker is asking about
\item \txx{reference}: the kind of entity the speaker is referring to
\end{itemize}
\item \txx{Syntax/semantics interface}: 
  \begin{quote}
    Constraints on how syntactic arguments are related to semantic
    ones, and on how semantic information is compiled from different
    parts of the sentence.
  \end{quote}
\end{itemize}

\myslide{Feature Geometry}
\avm{[\type*{expression}
syn & [\type*{syn-cat}
      head & [\type*{pos} ... ] \\
      val & [ spr & < ... > \\
               comps & < ... > ]]\\
      sem & [\type*{sem-cat}
      mode &  \{ prop, ques, dir, ref, none \} \\
      index &  \{ i, j, k, \ldots, s$_1$, s$_2$, \ldots \} \\
      restr & < ... > ]]
}



% \myslide{How the Pieces Fit Together}
% \begin{avm}\avmfont{\sc}
%   \< \eng{Kim}, \[syn & \[head & \[{\it noun}\\
%   agr & {\it 3sing}\]\\
%   val & \[spr & {< >}\\
%   comps & {< >}\]\]\\
%   sem & \[mode & ref\\
%   index & \ \ $i$\\
%   restr & \< \[reln & {\bf name}\\
%   % sit & \ \ $s$\\
%   name & Kim\\
%   named & \ \ $i$\] \> \]\] \>  %
% \end{avm}
% \myslide{How the Pieces Fit Together}
% \begin{avm}\avmfont{\sc}
% \< \eng{sleep} ,\ \[syn & \[head & {\it verb}\\
%                                      val & \[spr & \q< NP$_i$ \q>\\
%                                      comps & \q<  \q>\]\]\\
%                             sem & \[mode & prop\\
%                                     index & $s$\\
%                                     restr & \< \[reln & {\bf sleep}\\
%                                                  sit & \ \ $s$\\
%                                                  sleeper & \ \ $i$  \] \>\]\] \>
%                                         \end{avm}
% \myslide{The Pieces Together}

% \begin{avmtree}\avmfont{\sc}
%   \br{S}{
%     \br{\@{1} NP  \[ sem & \[ index & $i$ \] \] }{\lf{\eng{Kim}}}
%     \br{VP  \[ syn & \[ val & \[ spr & \< \@{1} \> \] \] \\
%       sem & \[ mode & prop\\
%       index & $s$\\
%       restr & \< \[reln & {\bf sleep}\\
%       sit & \ \ $s$\\
%       sleeper & \ \ $i$  \] \> \] \]    }{\lf{\eng{slept}}}}
% \end{avmtree}


\myslide{An Example}
\noindent\scalebox{0.7}{\begin{forest}
cat
[{S  \avm{[ sem & [ mode & prop \\ index & $s$\\
      restr & < [reln & {\bf name}\\
      name & \ \ Kim \\
      named & \ \ $i$ ],
      [reln & {\bf sleep}\\
      sit & \ \ $s$\\
      sleeper & \ \ $i$  ]  >  ] ]}} [{\avm{\1} NP  \avm{[ sem & [ mode & ref \\ index & $i$\\
      restr & < [reln & {\bf name}\\
      name & \ \ Kim \\
     named & \ \ $i$ ] >  ] ]}} [{\eng{Kim}}]]
[{\HD VP  \avm{[ syn & [ val & [ spr & < \1 > ] ] \\
      sem & [  mode & prop \\ index & $s$ \\ 
      restr & < [reln & {\bf sleep}\\
      sit & \ \ $s$\\
      sleeper & \ \ $i$  ] > ] ]}} [{\eng{slept}}]]]
\end{forest}}


\myslide{How to Share Semantic Information}
\MyLogo{List summation: $\oplus$ (technically concatenation)} 
\begin{itemize}
\item \txx{The Semantic Inheritance Principle}
  \begin{quote}
    In any headed phrase, the mother's \ft{mode} and
    \ft{index} are identical to those of the head daughter.    
  \end{quote}
\item \txx{The Semantic Compositionality Principle}
  \begin{quote}
    In any well-formed phrase structure, the mother's
    \ft{restr} value is the sum of the \ft{restr} values of
    the daughters.
  \end{quote}
\end{itemize}

\myslide{Where Is the Information?}
\MyLogo{}
\begin{itemize}
\item Words
  \begin{itemize}
  \item Contribute predications
  \item ‘Expose’ one index in those predications, for use by words or phrases
  \item Relate syntactic arguments to semantic arguments
  \end{itemize}
\item Rules
\begin{itemize}
\item Identify (link) feature structures across daughters
\item License trees that are subject to the semantic principles
  \begin{itemize}
  \item SIP: ‘passes up’ \ft{mode} and \ft{index} from head daughter
  \item SCP: ‘gathers up’ predications (\ft{restr} list) from all daughters
  \end{itemize}
\end{itemize}
\item The semantics is strictly compositional --- all of the meaning
  comes from the words, rules and principles.
\end{itemize}

\section{The Informal Binding Theory}


\myslide{Some Examples from Chapter 1}

\begin{exe}
\ix \eng{She likes herself.}
\ix \bad{She$_i$ likes her$_i$.}
\ix \eng{We gave presents to ourselves.}
\ix \bad{We gave presents to us.}
\ix \eng{We gave ourselves presents.}
\ix \bad{We gave us presents.}
\ix \bad{Leslie told us about us.}
\ix \eng{Leslie told us about ourselves.}
\ix \bad{Leslie told ourselves about us.}
\ix \bad{Leslie told ourselves about ourselves.}
\end{exe}

\myslide{Some Terminology}
\begin{itemize}
\item \txx{Binding}: The association between a pronoun
and an antecedent.
\item \txx{Anaphoric}: A term to describe an element (e.g.
a pronoun) that derives its interpretation from
some other expression in the discourse.
\item \txx{Antecedent}: The expression an anaphoric
expression derives its interpretation from.
\item \txx{Anaphora}: The relationship between an
anaphoric expression and its antecedent.
\end{itemize}

\myslide{The Chapter 1 Binding Theory Reformulated}
\begin{itemize}
\item Old formulation:
\begin{itemize}
\item A reflexive pronoun must be an argument of a verb that
has another preceding argument with the same reference.
\item A nonreflexive pronoun cannot appear as an argument of
a verb that has a preceding coreferential argument.
\end{itemize}
\item New formulation (version I):
\begin{itemize}
\item \txx{Principle A}: A reflexive pronoun must be
bound by a preceding argument of the same verb.
\item \txx{Principle B}: A nonreflexive pronoun may not
be bound by a preceding argument of the same verb.
\end{itemize}
\item The (opaque) names Principle A and B come from \citet{chomsky:1981}
\end{itemize}

\myslide{Some Challenges}
\begin{itemize}
\item Replace notions of \emp{bound} and \emp{preceding
argument of the same verb} by notions
definable in our theory.
\item Generalize the Binding Principles to get
better coverage.
\end{itemize}


\qaslide{Mapping the Concepts to Our Formalism}{
\begin{itemize}\addtolength{\itemsep}{-1.5ex}
\item[Q]  What would be a natural way to formalize
the notion of “bound” in our theory?
\answer{\item[A] Two expressions are bound if
they have the same \ft{index} value (“are
\emp{coindexed}”).}
\item[Q] Where in our theory do we have information
about a verb’s arguments?
\answer{\item[A] In the verb’s \ft{val} features:
  \ft{spr} (the subject) and \ft{comps} (the complements).}
\item[Q] What determines the linear ordering of a
verb’s arguments in a sentence?
\answer{\item[A] The interaction of the grammar
rules and the ordering of elements in the
\ft{comps} list:
\begin{itemize}
\item the Head-Specifier Rule puts the \ft{spr} \emp{before} the verb
\item the Head-Complement Rule puts the \ft{comps} \emp{after} the
  verb, \emp{in the order they are listed}
\end{itemize}}
\end{itemize}
}

\section{Formalizing Binding Theory}

\myslide{The Argument Realization Principle}
\begin{itemize}
\item For Binding Theory, we need a single list with both subject
and complements.
\item We introduce a feature \ft{arg-st} (\emp{argument structure}),
  with the following property (to be revised later): \\
\avm{[\type*{word}  
 syn & [ val & [ spr &  \tag{A} \\comps &  \tag{B} ] ] \\
 arg-st & \tag{A} $\oplus$ \tag{B} ]
}

\item In plain English: \emp{a word's \ft{arg-st} list is its \ft{spr}
    list followed by its \ft{comps} list}
  \begin{itemize}
  \item \eng{sent}: \ft{spr} $\langle$NP$\rangle$, \ft{comps}
    $\langle$NP, PP$\rangle$, so \ft{arg-st} $\langle$NP, NP, PP$\rangle$
  \end{itemize}
\item This is a constraint on the type \val{word}: every word has
  an \ft{arg-st} built this way
\end{itemize}

\myslide{Notes on ARG-ST}
\begin{itemize}
\item It’s neither in \ft{syn} nor \ft{sem}
  \begin{itemize}
  \item ``but rather serves to express certain relations at the
    interface of syntax and semantics''
    \citep[p.~207]{sag:wasow:bender:2003}
  \end{itemize}
\item It only appears on lexical items, i.e.\ words (not
  appropriate for type \val{phrase})
\item No principle stipulates identity
between \ft{arg-st}s: each word has its own, and nothing passes
them up the tree
\end{itemize}

\myslide{Two Bits of Technical Machinery}
\begin{itemize}
\item Definition of \txx{outranks}: \\
  If A precedes B on some \ft{arg-st} list,
  then A \txx{outranks} B.
  \begin{itemize}
  \item So the subject outranks all the complements, and the first
    complement outranks the second, \ldots
  \end{itemize}
\item Elements that must be \txx{anaphoric} --- that is, that
require an antecedent --- are lexically marked
[\ft{mode} \val{ana}]. 
\begin{itemize}
\item These include reflexive pronouns (\eng{myself, yourself,
    themselves, \ldots}) and reciprocals (\eng{each other, one another}).
\end{itemize}
\item Normal \txx{referential} NPs, including non-reflexive pronouns
  (\eng{I, me, them, \ldots}),  are marked [\ft{mode} \val{ref}].
\end{itemize}

\myslide{The Binding Principles}
\begin{itemize}
\item \txx{Principle A}: A [\ft{mode} \val{ana}] element must be
outranked by a coindexed element.
\begin{itemize}
\item In plain English: \emp{a reflexive needs an antecedent with the
    same index that comes earlier on an \ft{arg-st} list}
\end{itemize}
\item \txx{Principle B}: A [\ft{mode} \val{ref}] element must not
be outranked by a coindexed element.
\begin{itemize}
\item In plain English: \emp{an ordinary NP or pronoun must not have
    an antecedent with the same index earlier on an \ft{arg-st} list}
\end{itemize}
\item[Q] Compare these with version I: what has changed?
\item[A] ``bound'' is now \emp{coindexed} and ``preceding argument of
  the same verb'' is now \emp{outranks}, both defined in our theory
\end{itemize}

\myslide{Pronoun-Antecedent Agreement}
\begin{itemize}
\item The Binding Principles by themselves don’t block:
  \begin{exe}
    \ix \bad{I amused yourself.}
    \ix \bad{He amused themselves.}
    \ix \bad{She amused himself.}
  \end{exe}
\item Coindexed NPs refer to the same entity, and \ft{agr} features
generally correlate with properties of the referent.
\item The \txx{Anaphoric Agreement Principle} (AAP):
Coindexed NPs agree.
\end{itemize}

\section{Binding and PPs}

\qaslide{Binding in PPs}{
\begin{itemize}
\item What do the Binding Principles predict about the
following (so far)?
  \begin{exe}
    \ix \eng{I brought a book with me.}
    \ix \bad{I brought a book with myself.}
    \ix \bad{I mailed a book to me.}
    \ix \eng{I mailed a book to myself.}
  \end{exe}
\answer{\item[A] So far, the same thing for both pairs:
  \begin{itemize}
  \item the pronoun is only on the \ft{arg-st} of the preposition,
    where nothing outranks it
  \item so Principle A rules out \emp{both} reflexives, and Principle B
    allows \emp{both} plain pronouns
  \item That gets \eng{with} right and \eng{to} wrong: \emp{we need to tell
    the two prepositions apart}
  \end{itemize}}
\end{itemize}
}

\myslide{Two Types of Prepositions: the Intuition}
\begin{itemize}
\item \txx{Argument-marking}: Function like case markers in other languages, indicating the
roles of NP referents in the situation denoted by the verb.
\begin{itemize}
\item \eng{I mailed a book \ul{to} Kim}: \eng{to} just marks Kim as
  the recipient of the mailing; Japanese \mtciteform{ni}, Czech dative
\end{itemize}
\item \txx{Predicative}: Introduce their own predication.
\begin{itemize}
\item \eng{I brought a book \ul{with} me}: \eng{with} adds its own
  meaning --- a \textbf{with} relation between the book and me
\end{itemize}
\end{itemize}

\myslide{Two Types of Prepositions: a Formalization}
  \begin{itemize}
  \item Argument-marking prepositions share their
    objects' \ft{mode} and \ft{index} values.
    \begin{itemize}
    \item This is done with tagging in the lexical
      entries of such prepositions.
    \item These features are also shared with the PP
      node, by the Semantic Inheritance Principle.
    \item So the PP has the \emp{same index as the NP inside it}
    \end{itemize}
  \item Predicative prepositions introduce their own
\ft{mode} and \ft{index} values.
\begin{itemize}
\item So the PP has an index of its own (a situation), \emp{different
    from its object's}
\end{itemize}
\end{itemize}

\myslide{Argument-Marking Preposition \eng{to}}
\scalebox{0.9}{\avm{< \textit{to}, [ syn & [ head & prep \\
                     val &  [ spr & < > ] ] \\
                   arg-st & < NP [syn & [ head & [ case & acc ]  ]\\
                                  sem & [ mode & \1 \\ index & \2 ] ] >  \\
                   sem &  [ mode & \1 \\ index & \2 ] ]    >
}}
\hfill
\avm{\textsc{arg-st} = spr $\oplus$ comps
}

Therefore \ft{comps} here is a singleton list containing the NP, \ldots
\begin{itemize}\addtolength{\itemsep}{-1.5ex}
\item \avm{\1} and \avm{\2}: \eng{to} has \emp{the same \ft{mode} and
    \ft{index} as its object}; it has no predication of its own
\item By the SIP, the PP \eng{to myself} then also has \eng{myself}'s
  \ft{mode} and \ft{index}: it is PP$_i$[\ft{mode} \val{ana}]
\end{itemize}

\myslide{Redefining Rank}
 \begin{itemize}
\item If there is an \ft{arg-st} list on which A
precedes B, then A outranks B.
\item If a node is coindexed with its daughter, they
are of equal rank --- that is, they outrank the
same nodes and are outranked by the same
nodes.
\begin{itemize}
\item In plain English: \emp{for ranking, a PP headed by an
    argument-marking preposition counts as its object NP}
\item The verb's \ft{arg-st} only sees the PP, but the PP and the
  pronoun inside it share an index, so the pronoun inherits the PP's
  rank
\end{itemize}
\end{itemize}

\myslide{\eng{I sent a letter to myself}}

\begin{forest}
cat
[{S} [{\avm{\1} NP$_i$} [{\eng{I}}]]
[{VP \avm{[ spr & < \1 > ]}} [{V \avm{[ spr & < \1 > \\
             comps & < \2, \3 > \\ 
             arg-st & < \1, \2, \3 > ]}} [{\eng{sent}}]]
[{\avm{\2} NP$_j$} [{D} [{\eng{a}}]]
[{N$_j$} [{\eng{letter}}]]]
[{\avm{\3} PP$_i$} [{P$_i$} [{\eng{to}}]]
[{NP$_i$ \avm{[ mode & ana ]}} [{\eng{myself}}]]]]]
\end{forest}




\qaslide{The ARG-ST of \eng{sent} in This Tree}{
\begin{center}
  \avm{[  arg-st & < 
NP$_i$ [ mode & ref ],
NP$_j$ [ mode & ref ],
PP$_i$ [ mode & ana ]
> ]
}
\end{center}
\begin{itemize}
\item The PP is outranked by the first NP. (Why?)
  \answer{\\ \emp{Both are on \eng{sent}'s \ft{arg-st}, and the NP
      comes first}}
\item \eng{myself} has the same rank as the PP. (Why?)
  \answer{\\ \emp{\eng{to} is argument-marking, so the PP is coindexed
      with its daughter \eng{myself} (both $i$)}}
\item So, \eng{myself} is outranked by the first NP. (Why?)
  \answer{\\ \emp{Equal rank means outranked by the same things as
      the PP}}
\item Therefore, Principle A is satisfied:  A [\ft{mode} \val{ana}] element must be
outranked by a coindexed element.
  \answer{\\ \emp{\eng{myself} is [\ft{mode} \val{ana}], and it is
      outranked by NP$_i$ \eng{I}, which has the same index}}
\end{itemize}
}

\myslide{\eng{*I sent a letter to me}}
\begin{forest}
cat
[{S} [{\avm{\1} NP$_i$} [{\eng{I}}]]
[{VP \avm{[ spr & < \1 > ]}} [{V \avm{[ spr & < \1 > \\
             comps & < \2, \3 > \\ 
             arg-st & < \1, \2, \3 > ]}} [{\eng{sent}}]]
[{\avm{\2} NP$_j$} [{D} [{\eng{a}}]]
[{N$_j$} [{\eng{letter}}]]]
[{\avm{\3} PP$_i$} [{P$_i$} [{\eng{to}}]]
[{NP$_i$ \avm{[ mode & ref ]}} [{\eng{me}}]]]]]
\end{forest}

\myslide{The ARG-ST of \eng{sent} in This Tree}
\begin{center}
  \avm{[  arg-st & < 
NP$_i$ [ mode & ref ],
NP$_j$ [ mode & ref ],
PP$_i$ [ mode & ref ]
> ]
}
\end{center}
\begin{itemize}
\item The PP is outranked by the first NP (same \ft{arg-st}, NP first)
\item \eng{me} has the same rank as the PP (\eng{to} is
  argument-marking, so they are coindexed)
\item So, \eng{me} is outranked by the first NP.
\item Therefore, Principle B is violated: A [\ft{mode} \val{ref}] element must not
be outranked by a coindexed element.
  \begin{itemize}
  \item \eng{me} is [\ft{mode} \val{ref}], and NP$_i$ \eng{I}, with the
    same index, outranks it
  \end{itemize}
\end{itemize} 

\myslide{\eng{I brought a pencil with me}}
\MyLogo{Here \eng{I} does not outrank \eng{me}, so Principle B is satisfied}
\medskip

\scalebox{0.8}{\begin{forest}
cat
[{S} [{\avm{\1} NP$_i$} [{\eng{I}}]]
[{VP \avm{[ spr & < \1 > ]}} [{V \avm{[ spr & < \1 > \\
             comps & < \2, \3 > \\ 
             arg-st & < \1, \2, \3 > ]}} [{\eng{brought}}]]
[{\avm{\2} NP$_j$} [{D} [{\eng{a}}]]
[{N$_j$} [{\eng{pencil}}]]]
[{\avm{\3} PP$_k$} [{P$_k$ \avm{[ comps & < \4 > \\ 
                   arg-st & < \4 > ]}} [{\eng{with}}]]
[{\avm{\4} NP$_l$ \avm{[ mode & ref ]}} [{\eng{me}}]]]]]
\end{forest}}

\myslide{\eng{*I brought a pencil with myself}}
\MyLogo{Here \eng{I} does not outrank \eng{myself}, so Principle A is violated}
\medskip

\scalebox{0.8}{\begin{forest}
cat
[{S} [{\avm{\1} NP$_i$} [{\eng{I}}]]
[{VP \avm{[ spr & < \1 > ]}} [{V \avm{[ spr & < \1 > \\
             comps & < \2, \3 > \\ 
             arg-st & < \1, \2, \3 > ]}} [{\eng{brought}}]]
[{\avm{\2} NP$_j$} [{D} [{\eng{a}}]]
[{N$_j$} [{\eng{pencil}}]]]
[{\avm{\3} PP$_k$} [{P$_k$ \avm{[ comps & < \4 > \\ 
                   arg-st & < \4 > ]}} [{\eng{with}}]]
[{\avm{\4} NP$_l$ \avm{[ mode & ana ]}} [{\eng{myself}}]]]]]
\end{forest}}

\qaslide{The ARG-STs of \eng{brought} and \eng{with}}{
\MyLogo{}
\begin{center}
  \avm{[  arg-st & < 
NP$_i$ [ mode & ref ],
NP$_j$ [ mode & ref ],
PP$_k$ 
> ]}
\hspace{3em}
  \avm{[  arg-st & < 
NP$_l$ [ mode & ref / ana ]
> ]}
\end{center}
\begin{itemize}
\item The PP is outranked by the first NP. (Why?)
  \answer{\\ \emp{Both are on \eng{brought}'s \ft{arg-st}, and the NP comes
      first}}
\item But \eng{me}/\eng{myself} does \emp{not} have the same rank as
  the PP.  (Why?)
  \answer{\\ \emp{\eng{with} is predicative: PP$_k$ has its own index,
      so it is not coindexed with its daughter NP$_l$}}
\item \eng{me}/\eng{myself} is only on \eng{with}'s \ft{arg-st}, where
  nothing precedes it: \emp{nothing outranks it}
\item So Principle B is satisfied by \eng{me}, and Principle A is
  violated by \eng{myself}
\end{itemize}
}

\section{Imperatives}

\myslide{Imperatives}
\MyLogo{}
\begin{itemize}
\item Have the internal structure of a VP
  \begin{exe}
    \ix Leave!
    \ix Read a book!
    \ix Give the dog a treat!
    \ix Put the ice cream in the freezer!
  \end{exe}
\item Function as \txx{directives} (commands or requests)
\item Have the verb in base form: \eng{Be careful!} not \eng{*Are careful!}
\item Allow 2nd person reflexives, and no others
  \begin{exe}
    \ix \eng{Defend yourself!}
    \ix \bad{Defend myself/himself!}
  \end{exe}
\end{itemize}

\myslide{The Imperative Rule}

\avm{[\type*{phrase}  
   head & verb \\
   val & [ spr  & <  > ] \\
   sem & [ mode & dir \\ index & $s$ ] ]
              \ $\rightarrow$\ \ 
[\type*{phrase}  
   head & [\type*{verb} form & base ]\\
   val & [ spr  & < NP [ per & 2nd ]   > \\
            comps & < >  ] \\
   sem & [ index & $s$ ] ]
}


\begin{itemize}\addtolength{\itemsep}{-2ex}
\item Internal structure of a VP
\item Directive function
\item Base form
\item Only 2nd person reflexives
\item[Q] Note that this is not a headed rule. Why?
\item[A] It would violate the HFP and the SIP.
\end{itemize}

\qaslide{Imperative Example (Combining Constraints)}{
\noindent\begin{minipage}[t]{0.5\textwidth}\vspace{0pt}
\scalebox{0.8}{\begin{forest}
cat
[{S \avm{[ spr & < > ]}} [{VP \avm{[ spr & < \1 NP$_i$ [ per & 2nd \\ num & sg ] > ]}} [{V \avm{[ spr & < \1 > \\ comps & < \3 > \\ arg-st & < \1, \3 > ]}} [{\eng{vote}}]]
[{\avm{\3} PP$_i$} [{P$_i$} [{\eng{for}}]]
[{NP$_i$ \avm{[ mode & ana ]}} [{\eng{yourself}}]]]]]
\end{forest}}
\end{minipage}\hfill
\begin{minipage}[t]{0.47\textwidth}\vspace{0pt}
\begin{small}
\begin{itemize}
\item What's the \ft{spr} value on S?
  \answer{\\ \emp{$\langle\ \rangle$: the Imperative Rule says so}}
\item What's the \ft{spr} value on VP?
  \answer{\\ \emp{$\langle$NP[2nd]$\rangle$: \eng{vote}'s \ft{spr},
      passed up by the Valence Principle; the Imperative Rule
      adds [\ft{per} \val{2nd}]}}
\item What's the \ft{spr} value on V?
  \answer{\\ \emp{The same \avm{\1}: from \eng{vote}'s lexical entry}}
\item Which nodes have \ft{arg-st}?
  \answer{\\ \emp{Only the words: \eng{vote}, \eng{for}, \eng{yourself}}}
\item Which \ft{arg-st} matters for
the licensing of \eng{yourself}?
  \answer{\\ \emp{\eng{vote}'s: $\langle$NP$_i$, PP$_i$$\rangle$}}
\end{itemize}
\end{small}
\end{minipage}
}

\qaslide{\ft{arg-st} of \eng{vote}}{
\MyLogo{}
\begin{center}
  \avm{[  arg-st & < 
NP$_i$ [ mode & ref \\ per & 2nd \\ num & sg ],
PP$_i$ [ mode & ana ]
> ]
}
\end{center}
\begin{itemize}
\item Is Principle A satisfied?  How?
  \answer{\\ \emp{Yes: \eng{yourself} has the same rank as PP$_i$
      (\eng{for} is argument-marking), and PP$_i$ is outranked by the
      coindexed NP$_i$}}
\item Is Principle B satisfied?  How?
  \answer{\\ \emp{Yes: the only [\ft{mode} \val{ref}] element is the
      NP, and nothing outranks it}}
\item But NP$_i$ is never realised in the tree!  How can it be an
  antecedent?
  \answer{\\ \emp{Binding looks at \ft{arg-st}, not the tree: the
      subject is on \eng{vote}'s \ft{arg-st} even if unrealised.
      The AAP then makes \eng{yourself} 2nd person: \bad{Vote for myself!}}}
\end{itemize}
}

\myslide{Day 1 Revisited}
\begin{itemize}
\item Recall
  \begin{exe}
    \ix \eng{Screw yourself!}
    \ix \eng{Go screw yourself!}
    \ix \eng{Screw you!}
    \ix \bad{Go screw you!}
  \end{exe}
\item \eng{Screw NP!} has two analyses
  \begin{itemize}
  \item As an imperative
  \item As a truly subjectless fixed expression
  \end{itemize}
\item \eng{Go screw NP!} can only be analyzed as an imperative.
\end{itemize}


\myslide{Summary}
\MyLogo{Binding is much, much more complicated than this}
\begin{itemize}\addtolength{\itemsep}{-2ex}
\item Review of Chapter 1's informal binding
theory
\item Formalized Binding Theory
  \begin{itemize}
  \item \txx{Argument Realization Principle} (ARP)
    \\  \ft{arg-st} = \ft{spr}  $\oplus$ \ft{comps}
  \item \txx{Anaphoric Agreement Principle} (AAP)
   \\ Coindexed NPs agree.
 \item \txx{Principle A}\ \  A [\ft{mode} \val{ana}] element must be
outranked by a coindexed element.
\item \txx{Principle B}\ \  A [\ft{mode} \val{ref}] element must not
be outranked by a coindexed element.
  \end{itemize}
\item Binding and PPs
\item Imperatives
\end{itemize}


\section{Problem Set 6}

%%% Review exercise: semantic identity tables, after Emily Bender's
%%% 2026 in-class activity.  The tree is the one from lec-05.
\newcommand{\nd}[2]{{\setlength{\fboxsep}{1pt}\fcolorbox{#1}{white}{#2}}}
\newcommand{\matsem}{%
\scalebox{0.65}{\avm{[ mode & prop \\ index & $s_1$ \\
 restr & < [reln & {\bf the} \\ bv & $i$ ],
           [reln & {\bf cat} \\ inst & $i$ ],
           [reln & {\bf sleep} \\ sit & $s_1$ \\ sleeper & $i$ ],
           [reln & {\bf on} \\ sit & $s_2$ \\ located & $s_1$ \\ location & $j$ ],
           [reln & {\bf the} \\ bv & $j$ ],
           [reln & {\bf mat} \\ inst & $j$ ] > ]}}}
\newcommand{\mattree}{%
\scalebox{0.6}{\begin{forest}
cat
[{S} [{\avm{\1} NP$_i$} [{D$_i$} [{\eng{the}}]]
  [{Nom$_i$ \avm{[ spr & < \1 > ]}} [{\eng{cat}}]]]
[{VP \avm{[ spr & < \1 > ]}}
  [{\avm{\5} VP$_{s_1}$ \avm{[ spr & < \1 > ]}} [{\eng{slept}}]]
  [{PP$_{s_2}$ \avm{[ mod & < \5 > ]}} [{P$_{s_2}$ \avm{[ comps & < \6 > ]}} [{\eng{on}}]]
    [{\avm{\6} NP$_j$} [{D$_j$} [{\eng{the}}]] [{Nom$_j$} [{\eng{mat}}]]]]]]
\end{forest}}}

\myslide{P0: Semantic Identity Tables (Review)}
\MyLogo{After Emily Bender's 2026 in-class activity}
The \ft{sem} of \eng{The cat slept on the mat} (the tree from last week):
\begin{center}
\matsem
\end{center}
\begin{itemize}\addtolength{\itemsep}{-1ex}
\item Every index that appears twice in \ft{restr} got there by
  a \emp{chain of identities}: each link is required by a lexical
  entry, a grammar rule, or a principle
\item A \emp{semantic identity table} lists the links: what required
  the identity, and which two things it identifies
\item Let's trace \textbf{the}.\ft{bv} \& \textbf{cat}.\ft{inst} together
\end{itemize}

\qaslide{P0: \textbf{the}.\ft{bv} \& \textbf{cat}.\ft{inst}}{
\MyLogo{}
\begin{center}
\mattree
\end{center}
\answer{\par
\begin{footnotesize}
\begin{tabular}{@{}lll@{}}
\textbf{Grammar entity} & \textbf{First member} & \textbf{Second member} \\
\hline
lexical entry for \lex{cat} & \ft{inst} role of \textbf{cat} & \ft{index} of the word \lex{cat} ($i$) \\
lexical entry for \lex{cat} & \ft{index} of \lex{cat} & \ft{index} of the sole element of its \ft{spr} \\
Valence Principle & \lex{cat}'s \ft{spr} & Nom's \ft{spr} \avm{< \1 >} \\
Head-Specifier Rule & sole element of Nom's \ft{spr} \avm{\1} & the D node \lex{the} (and so its \ft{index}) \\
lexical entry for \lex{the} & \ft{index} of the word \lex{the} & \ft{bv} role of \textbf{the} \\
\end{tabular}
\end{footnotesize}}
}

\qaslide{P0: Three More Chains}{
\MyLogo{}
Now make a table for each of these pairs (links in order):
\begin{small}
\begin{itemize}\addtolength{\itemsep}{-1ex}
\item \textbf{mat}.\ft{inst} \& \textbf{on}.\ft{location}
  \answer{\\ \emp{entry for \lex{mat}: \ft{inst} = its \ft{index} ($j$)
      $\cdot$ SIP: Nom and NP get \lex{mat}'s \ft{index}
      $\cdot$ Head-Complement Rule: that NP is \eng{on}'s \ft{comps}
      element \avm{\6}
      $\cdot$ entry for \lex{on}: its complement's \ft{index} is the
      \ft{location}}}
\item \textbf{sleep}.\ft{sit} \& \textbf{on}.\ft{located}
  \answer{\\ \emp{entry for \lex{slept}: \ft{sit} = its \ft{index} ($s_1$)
      $\cdot$ SIP: the lower VP gets \lex{slept}'s \ft{index}
      $\cdot$ Head-Modifier Rule: that VP is \eng{on}'s \ft{mod}
      element \avm{\5}
      $\cdot$ entry for \lex{on}: its \ft{mod} element's \ft{index} is
      the \ft{located}}}
\item \textbf{cat}.\ft{inst} \& \textbf{sleep}.\ft{sleeper} \quad (last week's chain)
  \answer{\\ \emp{entry for \lex{slept}: \ft{sleeper} = its \ft{spr}'s
      \ft{index} $\cdot$ Valence Principle: upper VP has the same
      \ft{spr} $\cdot$ Head-Specifier Rule: \avm{\1} is that \ft{spr}
      $\cdot$ SIP: NP's \ft{index} is Nom's is \lex{cat}'s
      $\cdot$ entry for \lex{cat}: \ft{index} = \ft{inst}}}
\end{itemize}
\end{small}
}

\myslide{P1: English Possessives I}
\MyLogo{Based on \citet[Chapter 6, Problem 3]{sag:wasow:bender:2003}}

English uses {\it 's} to express possession, as in the following examples:

\begin{exe}
\ex  \eng{Leslie's coffee spilled.}
\ex  \eng{Jesse met the president of the university's cousin.}\label{s:i}
\ex  \bad{Jesse met the president's of the university cousin.}
\ex  \eng{Don't touch that plant growing by the trail's leaves.}\label{s:iv}
\ex \bad{Don't touch that plant's growing by the trail leaves.}
\ex  \eng{The person you were talking to's pants are torn.}\label{s:vi}
\ex  \bad{The person's you were talking to pants are torn.}\label{s:vii}
\end{exe}
(While examples (\ref{s:iv}) and (\ref{s:vi}) are a bit awkward, people do use such
sentences, and there is certainly nowhere else that the \eng{'s} could
be placed to improve them).  
\begin{itemize}
\item[A.] What is the generalization about where the \eng{'s} of
possession appears in English?
\end{itemize}
\newpage

% Consider noun \index{possessive}
% phrases like \eng{Kim's brother} and \eng{the president's salary}. 
One
traditional treatment of the possessive marker (\eng{'s})
% in English 
is to
claim it is a \index{case} case marker.
%-- which in our terms would mean it is a value of the feature \index{CASE!possessive}CASE.
In our terms this means that it indicates 
a particular value for the feature CASE (say, `poss' for `possessive')
on the word it attaches to.
% The principal objection to a treatment of \eng{'s}
% as a case marker is that, in conjunction with our analysis of case (in
% particular our assumption that \index{CASE} CASE is a HEAD feature),
% it makes the wrong 
% predictions about examples like the following:
%\begin{itemize}
%\item[B.] 
If we tried to formalize this traditional treatment of \eng{'s}, we
% would  need 
might posit
a rule along the following lines, 
based on the fact that possessive NPs appear in the same position as determiners:
\begin{itemize}
\item[] D \ $\rightarrow$\ \ \avm{NP [CASE & poss ]
}
\end{itemize}
Taken together with our assumption that \index{CASE}CASE is
a \index{head feature}HEAD feature,  such an analysis of \eng{'s} 
%as a case marker 
makes predictions about the grammaticality of (\ref{s:i})--(\ref{s:vii}).  
\begin{itemize}
\item[B.] Which of these sentences does it
predict should be grammatical, and why?

\end{itemize}

\myslide{P2: English Possessives II}
\MyLogo{Based on \citet[Chapter 6, Problem 4]{sag:wasow:bender:2003}}

An  alternative analysis of the possessive is to say that %it 
\eng{'s}
is a
\index{determiner (D)} determiner that builds a \index{determiner phrase
(DP)} determiner phrase (abbreviated DP),
via the \index{Head Specifier Rule@Head-Specifier Rule}Head-Specifier Rule. 
On this
analysis,  \eng{'s} selects for no complements, but it obligatorily
takes an NP \index{specifier} specifier. The word \eng{'s} thus has a lexical category
that is like an intransitive verb in valence.

This analysis is somewhat unintuitive, for two reasons:  first, it
requires that we have an independent lexical entry for \eng{'s}, which
seems more like a piece of a word, phonologically;  and second, it makes
the nonword \eng{'s} the head of its phrase!  However, this analysis
does a surprisingly good job of predicting the facts of English
 possessives, so we shall adopt it, at least for purposes of
this text.

\begin{itemize}
\item[A.] Ignoring semantics for the moment, give the lexical entry for
\eng{'s} assuming its analysis as a \index{determiner (D)} determiner, and draw 
%a phrasal SD (that
%is, 
a tree 
%description) 
for the NP \eng{Kim's brother}. 
%(Use abbreviations when
%necessary, for example: \hbox{NP[3rd, sg]} or
%\hbox{[SPR $\langle$NP[nom]$\rangle$]}.)
(The tree should show the value of \ft{head, spr} and \ft{comps} on every
node.  Use tags to show identities required by the grammar.
You may omit other features.)

\item[B.] Explain how your lexical entry gets the facts right in the 
% {\it Queen of England} examples above. 
following examples:
\begin{exe}
\ex \eng{The Queen of England's crown disappeared.}
\ex \bad{The Queen's of England crown disappeared.}
\end{exe}

\item[C.] How does this analysis handle recursion in possessives, 
for example, \eng{Robin's brother's wife}, or \eng{Robin's brother's wife's
parents}?  Provide at least one tree %description 
fragment
to illustrate your
explanation.
% (again you can abbreviate complex feature structures
% in your trees, 
%if the omitted details are clearly irrelevant).
(You may use abbreviations for node labels in the tree.)
\end{itemize}



% \myslide{P3: Prepositions}
% \MyLogo{Based on  Chapter 7, Problem 1, Sag, Wasow and Bender (2003)}

% For each of the following sentences, 
% \begin{itemize}
% \item[(a)] classify the underlined preposition into \txx{predicative}
%   or \txx{argument marking} and
% \item[(b)] justify by showing what reflexive and
%   nonreflexive coreferential 
%   pronouns can or cannot appear as the preposition's object.
% \end{itemize}
% \begin{itemize}\addtolength{\itemsep}{-1.5ex}
% \item[(i)] \eng{The dealer dealt an ace \ul{to} Bo.}
% \item[(ii)] \eng{The chemist held the sample away \ul{from} the flame.}
% \item[(iii)] \eng{Hiromi kept a flashlight \ul{beside} the bed.}
% \item[(iv)] \eng{We bought flowers \ul{for} you.}
% \item[(v)] \eng{The car has a scratch \ul{on} the fender.}
% \item[(vi)] \eng{The pronoun agrees \ul{with} its antecedent.}
% \end{itemize}

\myslide{P3: Imperative `Subjects'}
\MyLogo{Based on \citet[Chapter 7, Problem 2]{sag:wasow:bender:2003}}
\begin{itemize}
\item Some imperatives look like they have a subject:
  \begin{itemize}
  \item[(i)] \eng{\ul{You} get out of here!}
  \item[(ii)] \eng{\ul{Everybody} take out a sheet of paper!}
  \end{itemize}
\item But agreement is weird:
  \begin{itemize}
  \item[(iii)] 
    \eng{Everybody found 
      $^?$himself/$^*$yourself/themselves/$^*$myself a seat.}
  \item[(iv)]
    \eng{Everybody find 
      $^?$himself/yourself/$^*$themselves/$^*$myself a seat.}
  \end{itemize}
\item What minimal modification of the Imperative Rule would let these parse?
  \begin{itemize}
  \item Don't worry about the semantics for now (no one has solved that)
  \item We just want an imperative rule that will license the good examples and not the bad examples shown above.
    
  \end{itemize}


  
\end{itemize}



\myslide{Acknowledgments and References}

\begin{itemize}
\item Course design and slides borrow heavily from Emily Bender's course:
\textit{Linguistics 566: Introduction to Syntax for Computational Linguistics}
\\ \url{http://courses.washington.edu/ling566}
\item The worked examples, the imperative questions and the review
  exercise follow her Autumn 2026 slides (Binding \& Imperatives)
\item This lecture complements \citet[Chapter 7]{sag:wasow:bender:2003}
\end{itemize}


%\input{lec-05-questions}

\bibliographystyle{aclnat}
\bibliography{abb,mtg,nlp,ling}
\end{document}

%%% Local Variables: 
%%% coding: utf-8
%%% mode: latex
%%% TeX-PDF-mode: t
%%% TeX-engine: xetex
%%% End: 
